% ======================================================================
%  preambulo.tex — Estilo común para las tareas de CFD ICM3920
%  Uso: \input{preambulo} desde el archivo principal.
%  Compilar con: latexmk -pdf tarea-icm3920.tex
% ======================================================================

% ---------- Idioma y codificación ----------
\usepackage[utf8]{inputenc}
\usepackage[T1]{fontenc}
\usepackage[spanish,es-noshorthands]{babel}
\usepackage{lmodern}
\usepackage{microtype}
% ---------- Página ----------
\usepackage[a4paper,margin=2.5cm]{geometry}
\usepackage{fancyhdr}
\usepackage{parskip}

% ---------- Matemática y unidades ----------
\usepackage{amsmath,amssymb,bm}
\usepackage{siunitx}
\sisetup{
  detect-all,
  per-mode        = symbol,
  range-units     = single,
  separate-uncertainty = true,
  % Descomentar si el curso exige coma decimal:
  % output-decimal-marker = {,},
}
% Unidades de uso frecuente en CFD
\DeclareSIUnit{\poise}{P}
\DeclareSIUnit{\cell}{celda}

% ---------- Figuras y tablas ----------
\usepackage{graphicx}
\graphicspath{{figuras/}{figs/}}
\usepackage{float}
\usepackage{booktabs}
\usepackage{multirow}
\usepackage{longtable}
\usepackage{array}
\usepackage[font=small,labelfont=bf,width=0.9\textwidth]{caption}
\usepackage{tikz}
\usetikzlibrary{arrows.meta,calc}
\usepackage{subcaption}
\usepackage{enumitem}

% ---------- Color, cajas y código ----------
\usepackage{xcolor}
\definecolor{ofazul}{HTML}{1F4E79}
\definecolor{ofgris}{HTML}{F5F5F5}
\definecolor{ofverde}{HTML}{2E7D32}
\definecolor{ofrojo}{HTML}{B00020}
\definecolor{ofcomment}{HTML}{6A737D}

\usepackage{listings}
\lstset{
  basicstyle       = \ttfamily\footnotesize,
  backgroundcolor  = \color{ofgris},
  frame            = single,
  rulecolor        = \color{ofgris!60!black},
  breaklines       = true,
  breakatwhitespace= true,
  showstringspaces = false,
  columns          = fullflexible,
  keepspaces       = true,
  numbers          = none,
  captionpos       = b,
  extendedchars    = true,
  literate =
    {á}{{\'a}}1 {é}{{\'e}}1 {í}{{\'i}}1 {ó}{{\'o}}1 {ú}{{\'u}}1
    {Á}{{\'A}}1 {É}{{\'E}}1 {Í}{{\'I}}1 {Ó}{{\'O}}1 {Ú}{{\'U}}1
    {ñ}{{\~n}}1 {Ñ}{{\~N}}1 {ü}{{\"u}}1 {°}{{$^\circ$}}1
    % caracteres de árbol (salida de `tree`): se convierten a ASCII
    {│}{{|}}1 {├}{{|}}1 {└}{{|}}1 {─}{{-}}1 {—}{{-}}1
}

% Estilo para diccionarios de OpenFOAM (blockMeshDict, 0/U, fvSchemes, ...)
\lstdefinelanguage{OpenFOAM}{
  morekeywords = {FoamFile,version,format,class,object,location,dimensions,
    internalField,boundaryField,type,value,uniform,nonuniform,convertToMeters,
    vertices,blocks,edges,boundary,mergePatchPairs,hex,simpleGrading,arc,
    faces,neighbourPatch,inGroups,nFaces,startFace,
    ddtSchemes,gradSchemes,divSchemes,laplacianSchemes,interpolationSchemes,
    snGradSchemes,wallDist,solvers,SIMPLE,PISO,PIMPLE,relaxationFactors,
    application,startFrom,startTime,stopAt,endTime,deltaT,writeControl,
    writeInterval,writeFormat,writePrecision,runTimeModifiable,adjustTimeStep,
    maxCo,functions,simulationType,RAS,model,turbulence,printCoeffs},
  morekeywords = [2]{fixedValue,zeroGradient,noSlip,slip,empty,symmetry,cyclic,
    inletOutlet,pressureInletOutletVelocity,calculated,
    kqRWallFunction,epsilonWallFunction,omegaWallFunction,nutkWallFunction,
    totalPressure,fixedFluxPressure,freestream,movingWallVelocity},
  sensitive    = true,
  morecomment  = [l]{//},
  morecomment  = [s]{/*}{*/},
  morestring   = [b]",
}
\lstdefinestyle{foam}{
  language     = OpenFOAM,
  keywordstyle = \color{ofazul}\bfseries,
  keywordstyle = [2]\color{ofverde},
  commentstyle = \color{ofcomment}\itshape,
  stringstyle  = \color{ofrojo},
}
\lstdefinestyle{terminal}{
  language     = bash,
  basicstyle   = \ttfamily\footnotesize\color{black},
  keywordstyle = \color{ofazul}\bfseries,
  commentstyle = \color{ofcomment}\itshape,
  morekeywords = {blockMesh,checkMesh,snappyHexMesh,surfaceFeatureExtract,
    decomposePar,reconstructPar,simpleFoam,icoFoam,pisoFoam,pimpleFoam,
    interFoam,potentialFoam,foamToVTK,postProcess,paraFoam,setFields,
    mpirun,foamCleanTutorials,transformPoints,refineMesh},
}

% Nombres en español para los listados
\renewcommand{\lstlistingname}{Listado}
\renewcommand{\lstlistlistingname}{Índice de listados}

% Atajos: \begin{foamcode}{ruta/archivo} ... \end{foamcode}
\lstnewenvironment{foamcode}[1][]
  {\lstset{style=foam,#1}}
  {}
\lstnewenvironment{terminal}[1][]
  {\lstset{style=terminal,#1}}
  {}

% ---------- Referencias cruzadas ----------
\usepackage[hidelinks]{hyperref}
\hypersetup{
  colorlinks = true,
  linkcolor  = ofazul,
  citecolor  = ofverde,
  urlcolor   = ofazul,
}
\usepackage[spanish,capitalise,noabbrev]{cleveref}
\crefname{figure}{Figura}{Figuras}
\crefname{table}{Tabla}{Tablas}
\crefname{equation}{Ec.}{Ecs.}
\crefname{section}{Sección}{Secciones}
\crefname{lstlisting}{Listado}{Listados}

% ---------- Macros propias ----------
\newcommand{\OF}{\textsc{OpenFOAM}}
\newcommand{\PV}{\textsc{ParaView}}
\newcommand{\archivo}[1]{\texttt{#1}}          % rutas y diccionarios
\newcommand{\comando}[1]{\texttt{#1}}          % utilidades del terminal

\newcommand{\Rey}{\ensuremath{\mathrm{Re}}}
\newcommand{\Cou}{\ensuremath{\mathrm{Co}}}
\newcommand{\Mach}{\ensuremath{\mathrm{Ma}}}
\newcommand{\Pran}{\ensuremath{\mathrm{Pr}}}
\newcommand{\Cd}{\ensuremath{C_D}}
\newcommand{\Cl}{\ensuremath{C_L}}
\newcommand{\St}{\ensuremath{\mathrm{St}}}
\newcommand{\yplus}{\ensuremath{y^{+}}}
\newcommand{\vect}[1]{\ensuremath{\bm{#1}}}
\newcommand{\dpar}[2]{\ensuremath{\dfrac{\partial #1}{\partial #2}}}

% Marca visible para lo que falta completar (comentar antes de entregar)
\newcommand{\pend}[1]{\textcolor{ofrojo}{[\textbf{PENDIENTE:} #1]}}

% ---------- Datos de la tarea (editar en el archivo principal) ----------
\newcommand{\curso}{ICM3920 — Fluido Dinamica Computacional}
\newcommand{\tareanum}{1}
\newcommand{\tareatitulo}{Título de la tarea}
\newcommand{\alumno}{Piter}
\newcommand{\profesor}{}
\newcommand{\fechaentrega}{\today}

% ---------- Encabezado y pie ----------
\pagestyle{fancy}
\fancyhf{}
\fancyhead[L]{\small\curso}
\fancyhead[R]{\small Tarea \tareanum}
\fancyfoot[C]{\thepage}
\renewcommand{\headrulewidth}{0.4pt}
